% ============================================================
%   Appendix E  --  Mathematical Formulation
% ============================================================

\chapter{Mathematical Formulation}\label{app:math}

For readability, several quantities in the main text are presented through figures and tables
rather than equations. This appendix reproduces their exact closed forms for reference, with every
symbol defined. In each case the figure cited in the main text is the primary, intuition-carrying
presentation, and the equations below are its algebraic counterpart.

\section{Risk-level classification (adaptive shield)}\label{app:math-risk}

The frontal-distance criterion of the adaptive shield (Section~\ref{sec:design-adaptive-shield},
Figure~\ref{fig:risk-levels} and Table~\ref{tab:risk-params}) maps the minimum of the dynamic-only
LIDAR front sector \(d^{\mathrm{dyn}}_t\in[0,1]\) to a discrete risk level:
\begin{equation}
  \mathrm{level}^{\mathrm{frontal}}_t \;=\;
     \begin{cases}
       \mathrm{safe}     & \text{if } d^{\mathrm{dyn}}_t > 0.50,\\
       \mathrm{warning}  & \text{if } 0.20 < d^{\mathrm{dyn}}_t \leq 0.50,\\
       \mathrm{critical} & \text{if } d^{\mathrm{dyn}}_t \leq 0.20,
     \end{cases}
  \label{eq:shield-risk-front}
\end{equation}
where a higher \(d^{\mathrm{dyn}}_t\) denotes a more distant in-path obstacle (\(1\) is a clear
lane). The lateral-position criterion applies the analogous thresholds of
Table~\ref{tab:risk-params} to the absolute lateral offset \(|\ell_t|\) and the absolute heading
error \(|\psi_t|\); the overall risk level is the worst (most severe) of the two criteria.

\section{Lane-keep gate (reward shaper)}\label{app:math-gate}

The lane-keep gate of Section~\ref{sec:design-reward}, visualised in
Figure~\ref{fig:lane-keep-gate}, multiplies the forward-progress bonuses by
\begin{equation}
  \mu_t \;=\;
  \underbrace{\max(1 - |\ell_t|,\, 0)}_{\text{centering factor}}
  \;\times\;
  \underbrace{\max\!\left(1 - \frac{|\psi_t|}{20^\circ},\, 0\right)}_{\text{heading factor}}
  \;\in\;[0,1],
  \label{eq:lane-keep-gate}
\end{equation}
where \(\ell_t\) is the normalised lateral offset (\(0\) at the lane centre, \(\pm1\) at the edges)
and \(\psi_t\) is the heading error in degrees. The gate equals \(1\) only when the vehicle is both
centred and aligned, and \(0\) as soon as \(|\ell_t|\geq1\) or \(|\psi_t|\geq20^\circ\). During a
permitted lane change (\(\ell_t>0.5\) with CARLA reporting the change allowed) \(\mu_t\) is floored
at \(0.3\).

\section{Shaped reward (reward shaper)}\label{app:math-reward}

The total shaped reward of Section~\ref{sec:design-reward}, summarised in
Figure~\ref{fig:reward-overview} and Table~\ref{tab:reward-terms}, is the sum of the sparse base
reward and the dense bonus and penalty terms:
\begin{equation}\label{eq:shaped-reward}
  \begin{split}
  \tilde{r}_t \;=\;{} & r^{\mathrm{base}}_t
    + r^{\mathrm{prog}}_t + r^{\mathrm{accel}}_t + r^{\mathrm{speed}}_t
    + r^{\mathrm{lane}}_t + r^{\mathrm{head}}_t + r^{\mathrm{mile}}_t \\
    & - p^{\mathrm{smooth}}_t - p^{\mathrm{inv}}_t - p^{\mathrm{solid}}_t
      - p^{\mathrm{edge}}_t - p^{\mathrm{wrong\text{-}h}}_t
      - p^{\mathrm{idle}}_t - p^{\mathrm{drift}}_t
      - p^{\mathrm{lc}}_t - p^{\mathrm{shield}}_t.
  \end{split}
\end{equation}
Each symbol is defined in Table~\ref{tab:reward-terms}; the default weights are listed in
Appendix~\ref{app:hyperparameters}, Table~\ref{tab:hp-reward}.

\section{Shield permissiveness (LLM meta-controller)}\label{app:math-permissiveness}

The single permissiveness scalar \(p\in[0,1]\) of Section~\ref{sec:dep-llm}, illustrated in
Figure~\ref{fig:permissiveness-ramp}, sets every tunable threshold \(\tau_i\) of the adaptive
shield by linear interpolation between a strict and a permissive endpoint:
\begin{equation}\label{eq:permissiveness}
  \tau_i(p) \;=\; \tau_i^{\mathrm{strict}} \;+\; (1-p)\,\bigl(\tau_i^{\mathrm{perm}} - \tau_i^{\mathrm{strict}}\bigr),
\end{equation}
where \(\tau_i^{\mathrm{strict}}\) and \(\tau_i^{\mathrm{perm}}\) are the strict and permissive
endpoints of threshold \(i\). Thus \(p=1\) recovers the production shield (every
\(\tau_i=\tau_i^{\mathrm{strict}}\)) and \(p=0\) the maximally weaned shield (every
\(\tau_i=\tau_i^{\mathrm{perm}}\)).
